% Part: incompleteness
% Chapter: arithmetization-syntax
% Section: proofs-in-lk

\documentclass[../../../include/open-logic-section]{subfiles}

\begin{document}

\olfileid[tr]{inc}{art}{plk}
\olsection{$\Log{LK}$ Sisteminde \usetoken{P}{derivation}}

\begin{explain}
!!{derivation}s{}i aritmetikleştirmek için !!{derivation}s{}i sayılarla temsil
etmemiz gerekir. Her çıkarımın ayrıca bir etiket taşıdığı sekant ağaçları olan
!!{derivation}s için en açık yaklaşım özyinelemeli bir temsildir: Bir
!!{derivation}{}i, bileşenleri son sekant, etiket ve son çıkarımın öncüllerine
götüren alt-!!{derivation}s{}in temsilleri olan bir demetle temsil ederiz.
\end{explain}

\begin{defn}
$\Gamma$ sonlu bir !!{sentence}s dizisi ve $\Gamma =
\tuple{!A_1, \dots, !A_n}$ ise $\Gn{\Gamma} = \tuple{\Gn{!A_1},
  \dots, \Gn{!A_n}}$ olur.

$\Gamma \Sequent \Delta$ bir sekant ise $\Gamma \Sequent \Delta$'nın Gödel numarası
\[
\Gn{\Gamma \Sequent \Delta} = \tuple{\Gn{\Gamma}, \Gn{\Delta}}
\]
olarak tanımlanır.

$\pi$, $\Log{LK}$ sisteminde !!a{derivation} ise $\Gn{\pi}$ aşağıdaki gibi
tanımlanır:
\begin{enumerate}
\item $\pi$ yalnızca $\Gamma \Sequent \Delta$ başlangıç sekantından
  oluşuyorsa $\Gn{\pi}$
  \[
    \tuple{0, \Gn{\Gamma \Sequent \Delta}}
  \]
  olur.
\item $\pi$, bir ya da iki öncüllü bir çıkarımla bitiyorsa, sonucu
  $\Gamma \Sequent \Delta$ ise ve $\pi_1$ ile $\pi_2$ son çıkarımın
  öncülleriyle biten dolaysız alt ispatlar ise $\Gn{\pi}$ sırasıyla
  \begin{align*}
    & \tuple{1, \Gn{\pi_1}, \Gn{\Gamma \Sequent \Delta}, k} \text{ ya da}\\
    & \tuple{2, \Gn{\pi_1}, \Gn{\pi_2}, \Gn{\Gamma \Sequent \Delta}, k}
  \end{align*}
  olur. Burada $k$, son çıkarımda kullanılan kurala göre aşağıdaki tabloyla
  verilir:

  \begin{tabular}{lccccccc}
    \text{Kural:} & \LeftR{\Weakening} & \RightR{\Weakening} &
    \LeftR{\Contraction} & \RightR{\Contraction} &
    \LeftR{\Exchange} & \RightR{\Exchange} \\
    $k$: & 1 & 2 & 3 & 4 & 5 & 6 \\[2ex]
    \text{Kural:} & \LeftR{\lnot} & \RightR{\lnot} &
    \LeftR{\land} & \RightR{\land} &
    \LeftR{\lor} & \RightR{\lor} \\
    $k$: & 7 & 8 & 9 & 10 & 11 & 12 \\[2ex]
    \text{Kural:} & \LeftR{\lif} & \RightR{\lif} &
    \LeftR{\lforall} & \RightR{\lforall} &
    \LeftR{\lexists} & \RightR{\lexists} \\
    $k$: & 13 & 14 & 15 & 16 & 17 & 18 \\[2ex]
    \text{Kural:} & \Cut & = \\
    $k$: & 19 & 20
  \end{tabular}
\end{enumerate}
\end{defn}

\begin{ex}
  Şu çok basit !!{derivation}{}i ele alın:
  \begin{prooftree}
    \Axiom$!A \fCenter !A$
    \RightLabel{\LeftR{\land}}
    \UnaryInf$!A \land !B \fCenter !A$
    \RightLabel{\RightR{\lif}}
    \UnaryInf$\fCenter (!A \land !B) \lif !A$
  \end{prooftree}
  Başlangıç sekantının Gödel numarası $p_0 = \tuple{0,
    \Gn{!A \Sequent !A}}$ olurdu. $\LeftR{\land}$ sonucuyla biten
  !!{derivation}{}in Gödel numarası $p_1 =
    \tuple{1, p_0, \Gn{!A \land !B \Sequent !A}, 9}$ olurdu
  ($\LeftR{\land}$ bir öncüllü olduğu için $1$, sonucun yani
  $!A \land !B \Sequent !A$'nın Gödel numarası ve $\LeftR{\land}$'ı
  kodlayan sayı olduğu için $9$). Böylece bütün !!{derivation}{}in Gödel
  numarası $\tuple{1, p_1, \Gn{\Sequent (!A \land !B) \lif !A)}, 14}$,
  yani
  \[
  \tuple{1, \tuple{1, \tuple{0, \Gn{!A \Sequent !A}}, \Gn{!A \land !B \Sequent !A}, 9},
    \Gn{\Sequent (!A \land !B) \lif !A}, 14}
  \]
  olur.
\end{ex}

\begin{explain}
!!{derivation}s için bir temsil belirlediğimize göre bu !!{derivation}s{}i
ilkel özyinelemeli biçimde işleyebildiğimizi ve temel özellikleriyle
bağıntılarını da bu biçimde ifade edebildiğimizi göstermeliyiz. Bazı işlemler
basittir: Örneğin !!a{derivation}{}in $p$ Gödel numarası verildiğinde
$\fn{EndSequent}(p) = (p)_{(p)_0+1}$ bize son sekantının Gödel numarasını,
$\fn{LastRule}(p) = (p)_{(p)_0+2}$ ise son kuralının kodunu verir.
\[
  \len{s} = 2 \land \bforall{i<\len{(s)_0} + \len{(s)_1}}{\fn{Sent}(((s)_0 \concat (s)_1)_i)}
\]
ile tanımlanan $\fn{Sequent}(s)$ özelliği, ancak ve ancak $s$ !!{sentence}s{}den
oluşan bir sekantın Gödel numarasıysa geçerlidir. Bazı işlemler ise çok daha
zordur. En azından bunların nasıl yapılacağını ana hatlarıyla göstereceğiz.
Amaç, ``$\pi$, $\Gamma$'dan $!A$'nın !!a{derivation}{}idir'' bağıntısının
$\pi$ ile $!A$'nın Gödel numaraları üzerinde ilkel özyinelemeli bir bağıntı
olduğunu göstermektir.
\end{explain}

\begin{prop}\ollabel{prop:followsby}
  Gödel numarası $p$ olan $\pi$ !!{derivation}{}indeki son çıkarımın doğru
  olması durumunda geçerli olan $\fn{Correct}(p)$ özelliği ilkel
  özyinelemelidir.
\end{prop}

\begin{proof}
  $\Gamma \Sequent \Delta$, ya $\Gamma \Sequent \Delta$ sekantı
  $!A \Sequent !A$ olacak biçimde !!a{sentence}~$!A$ varsa ya da
  $\Gamma \Sequent \Delta$ sekantı $\emptyset \Sequent \eq[t][t]$ olacak
  biçimde bir $t$ terimi varsa bir başlangıç sekantıdır. Gödel numaralarıyla
  $\fn{InitSeq}(s)$ ancak ve ancak
  \begin{align*}
  \bexists{x < s}{} (\fn{Sent}(x) & \land
  s = \tuple{\tuple{x},\tuple{x}})
  \lor {}\\
  \bexists{t<s}{} (\fn{Term}(t) & \land
  s = \tuple{0, \tuple{\Gn{{\eq}(} \concat t \concat \Gn{,} \concat t \concat \Gn{)}}})
  \end{align*}
  geçerliyse geçerlidir.

  Ayrıca her $R$ çıkarım kuralı için $\fn{FollowsBy}_R(p)$ bağıntısının
  ilkel özyinelemeli olduğunu göstermeliyiz. Burada
  $\fn{FollowsBy}_R(p)$, ancak ve ancak $p$, bir $\pi$ !!{derivation}{}inin
  Gödel numarasıysa ve $\pi$'nin son sekantı, $\pi$'nin dolaysız
  alt-!!{derivation}s{}inden $R$'nin doğru uygulanmasıyla çıkıyorsa geçerlidir.

  Basit bir durum $\RightR{\land}$ kuralıdır. $\pi$ doğru bir
  $\RightR{\land}$ çıkarımıyla bitiyorsa şu biçimdedir:
  \begin{prooftree}
    \AxiomC{}
    \RightLabel{$\pi_1$}
    \Deduce$\Gamma \fCenter \Delta, !A$

    \AxiomC{}
    \RightLabel{$\pi_2$}
    \Deduce$\Gamma \fCenter \Delta, !B$

    \RightLabel{\RightR\land}
    \BinaryInf$\Gamma \fCenter \Delta, !A \land !B$
  \end{prooftree}
  Dolayısıyla $\pi$ !!{derivation}{}indeki son çıkarımın
  $\RightR{\land}$ kuralının doğru bir uygulaması olması, ancak ve ancak
  $\pi_1$'in son sekantı $\Gamma \Sequent \Delta, !A$, $\pi_2$'nin son
  sekantı $\Gamma \Sequent \Delta, !B$ ve $\pi$'nin son sekantı
  $\Gamma \Sequent \Delta, !A \land !B$ olacak biçimde $\Gamma$ ve
  $\Delta$ !!{sentence}s dizileri ile iki $!A$ ve $!B$ !!{sentence}{}si varsa
  geçerlidir. Bunu yalnızca Gödel numaralarına çevirmemiz gerekir.
  $s = \Gn{\Gamma \Sequent \Delta}$ ise $(s)_0 = \Gn{\Gamma}$ ve
  $(s)_1 = \Gn{\Delta}$ olur. Öyleyse
  $\fn{FollowsBy}_{\RightR{\land}}(p)$ ancak ve ancak
\begin{align*}
& \bexists{g < p}{\bexists{d < p}{\bexists{a < p}{\bexists{b < p}{\quad}}}} \\
& \qquad \fn{EndSequent}(p) =
  \tuple{g, d \concat
    \tuple{\Gn{(} \concat a \concat \Gn{\land} \concat b \concat \Gn{)}}} \land {} \\
& \qquad \fn{EndSequent}((p)_1) = \tuple{g, d \concat \tuple{a}} \land {}\\
& \qquad \fn{EndSequent}((p)_2) = \tuple{g, d \concat \tuple{b}} \land {}\\
& \qquad (p)_0 = 2 \land \fn{LastRule}(p) = 10
\end{align*}
geçerliyse geçerlidir. Satırlar sırasıyla şunları ifade eder: Gödel numarası
$g$ olan bir $\Gamma$ dizisi, Gödel numarası $d$ olan bir $\Delta$ dizisi,
Gödel numarası $a$ olan !!a{formula}~$!A$ ve Gödel numarası $b$ olan
!!a{formula}~$!B$ vardır; ``$\pi$'nin son sekantı
$\Gamma \Sequent \Delta, !A \land !B$'dir'', ``$\pi_1$'in son sekantı
$\Gamma \Sequent \Delta, !A$'dır'', ``$\pi_2$'nin son sekantı
$\Gamma \Sequent \Delta, !B$'dir'' ve ``$\pi$'nin iki dolaysız alt
alt türetimi vardır ve son çıkarım kuralı (numarası~$10$ olan)
$\RightR\land$ kuralıdır.''

$\pi$'deki son çıkarımın $\RightR\lexists$ kuralının doğru bir uygulaması
olması, ancak ve ancak $\pi$'nin son sekantı
$\Gamma \Sequent \Delta, \lexists[x][!A]$ ve $\pi_1$'in son sekantı
$\Gamma \Sequent \Delta, \Subst{!A}{t}{x}$ olacak biçimde $\Gamma$ ile
$\Delta$ dizileri, !!a{formula}~$!A$, bir $x$ değişkeni ve bir $t$ terimi
varsa geçerlidir. Dolayısıyla Gödel numaralarıyla
$\fn{FollowsBy}_{\RightR\lexists}(p)$ ancak ve ancak
\begin{align*}
  & \bexists{g<p}{\bexists{d<p}{\bexists{a<p}{\bexists{x<p}{\bexists{t<p}{\quad}}}}}\\
  & \qquad \fn{EndSequent}(p) =
  \tuple{
    g,
    d \concat \tuple{\Gn{\lexists} \concat x \concat a}
  } \land {}\\
  & \qquad \fn{EndSequent}((p)_1) =
  \tuple{
    g,
    d \concat \tuple{\fn{Subst}(a, t, x)}
  } \land {}\\
  & \qquad (p)_0 = 1 \land \fn{LastRule}(p) = 18
\end{align*}
geçerliyse geçerlidir. Ardından $\fn{Correct}(p)$'yi şöyle tanımlarız:
\begin{multline*}
  \fn{Sequent}(\fn{EndSequent}(p)) \land {}\\
  [(\fn{LastRule}(p) = 1 \land
  \fn{FollowsBy}_{\LeftR\Weakening}(p)) \lor \dots \lor {}\\
  (\fn{LastRule}(p) = 20 \land \fn{FollowsBy}_{\eq}(p)) \lor {}\\
  (p)_0 = 0 \land \fn{InitSeq}(\fn{EndSequent}(p))]
\end{multline*}
İlk satır $p$'nin son sekantının gerçekten !!{sentence}s{}den oluşan bir
sekant olmasını sağlar. Son satır ise $p$'nin yalnızca bir başlangıç sekantı
olduğu durumu kapsar.
\end{proof}

\begin{prob}
  Aşağıdaki özellikleri \olref[inc][art][plk]{prop:followsby}'deki gibi
  tanımlayın:
  \begin{enumerate}
  \item $\fn{FollowsBy}_{\Cut}(p)$,
  \item $\fn{FollowsBy}_{\LeftR\lif}(p)$,
  \item $\fn{FollowsBy}_{\eq}(p)$,
  \item $\fn{FollowsBy}_{\RightR{\lforall}}(p)$.
  \end{enumerate}
  Sonuncusu için Gödel numarası $p$ olan !!{derivation}{}in son çıkarımının
  özdeğişken koşulunu sağlayıp sağlamadığını, yani
  $\RightR{\lforall}$ kuralının $a$ özdeğişkeninin son sekantta geçmediğini
  ilkel özyinelemeli biçimde sınayabildiğinizi de göstermelisiniz.
\end{prob}

\begin{prop}
  \ollabel{prop:deriv}
  $p$'nin doğru bir $\pi$ !!{derivation}{}inin Gödel numarası olması durumunda
  geçerli olan $\fn{Deriv}(p)$ bağıntısı ilkel özyinelemelidir.
\end{prop}

\begin{proof}
  !!^a{derivation}~$\pi$, ancak ve ancak çıkarımlarının her biri bir kuralın
  doğru uygulamasıysa, başka bir deyişle alt-!!{derivation}s{}inin her biri
  doğru bir çıkarımla bitiyorsa doğrudur. Dolayısıyla $\fn{Deriv}(p)$ ancak
  ve ancak
  \[
  \bforall{i<\len{\fn{SubtreeSeq}(p)}}{\fn{Correct}((\fn{SubtreeSeq}(p))_i)}
  \]
  geçerliyse geçerlidir.
\end{proof}


\begin{prop}
$\Gamma$'nın ilkel özyinelemeli bir !!{sentence}s kümesi olduğunu varsayın.
O zaman ``$x$, sonlu bir $\Gamma_0 \subseteq \Gamma$ için
$\Gamma_0 \Sequent !A$'nın !!a{derivation}{}i olan $\pi$'nin kodudur ve $y$,
$!A$'nın Gödel numarasıdır'' ifadesini veren $\Prf[\Gamma](x, y)$ bağıntısı
ilkel özyinelemelidir.
\end{prop}

\begin{proof}
``$y \in \Gamma$'' ifadesinin ilkel özyinelemeli $R_\Gamma(y)$ yüklemiyle
verildiğini varsayın. Ancak ve ancak $y$, bir $!A$ tümcesinin Gödel
numarasıysa ve $x$, son sekantı $\Gamma_0 \Sequent !A$ olan bir
$\Log{LK}$-!!{derivation}{}inin koduysa geçerli olan $\Prf[\Gamma](x, y)$
bağıntısının ilkel özyinelemeli olduğunu göstermeliyiz.

Önceki önermeye göre $x$'in $\Log{LK}$ sisteminde doğru bir $\pi$
!!{derivation}{}inin kodu olması durumunda geçerli olan $\fn{Deriv}(x)$
özelliği ilkel özyinelemelidir. $x$ böyle bir kodsa
$\fn{EndSequent}(x)$, $\pi$'nin son sekantının kodudur; dolayısıyla
$(\fn{EndSequent}(x))_0$ son sekantın sol tarafının, 
$(\fn{EndSequent}(x))_1$ ise sağ tarafının kodudur. Öyleyse
``$\pi$'nin son sekantının sağ tarafı $!A$'dır'' ifadesini
$\len{(\fn{EndSequent}(x))_1} = 1 \land ((\fn{EndSequent}(x))_1)_0 = y$
olarak yazabiliriz. $\pi$'nin son sekantının sol tarafı zaten sonludur;
yalnızca içindeki her tümcenin $\Gamma$'da olduğunu ifade etmemiz gerekir.
Böylece $\Prf[\Gamma](x, y)$'yi
\begin{align*}
\Prf[\Gamma](x, y) \defiff {}&
\fn{Deriv}(x) \land {} \\
& \bforall{i <
  \len{(\fn{EndSequent}(x))_0}}{R_\Gamma(((\fn{EndSequent}(x))_0)_i)} \land {}\\
& \len{(\fn{EndSequent}(x))_1} = 1 \land ((\fn{EndSequent}(x))_1)_0 = y
\end{align*}
olarak tanımlayabiliriz.
\end{proof}

\end{document}
